\section{Recoverability: what determines whether a shortcut can be repaired}
\label{sec:recoverability}

Results~\ref{sec:rq4}--\ref{sec:repair} are not three isolated outcomes; together they characterize \emph{when} a
verified harmful shortcut is repairable by a \emph{deployable, target-unlabeled} operator. We make the boundary
explicit. The deployable repairs we use are \textbf{target-$X$-only batch-statistic affine operators} $R(z)=Az+b$
that apply the \emph{same} map to every sample, with $A,b$ measurable from the target batch's moments
$\{\hat\mu_T,\hat\Sigma_T\}$ and source references (E4 is the translation case, E4b the linear case). Within this
class, a clean boundary holds.

\begin{quote}
\emph{Characterization (within the batch-affine operator and affine-in-$z$ corruption classes).} A corruption
$C:z_i\mapsto z_i+p_i$ is invertible by a target-$X$-only batch-affine operator iff $C$ is an affine function of
$z$ --- a batch-constant offset (undone by a translation) or a fixed \emph{full-rank} linear map (undone by a
matrix). An additive per-sample perturbation \emph{independent} of $z$ is not invertible, because the per-sample
realization is not present in the batch statistics.
\end{quote}

This places the empirical results into non-exclusive \emph{obstruction classes} (Table~\ref{tab:recover}). A
\textbf{first-moment deterministic} offset is affine-in-$z$ and, when task-separable and identifiable, is
recoverable (\textbf{R0}; E4, Result~\ref{sec:repair}); the recovery is narrow precisely because the offset is
conflated with the genuine domain-mean shift. A \textbf{per-sample stochastic} perturbation independent of $z$ is
\emph{not} invertible by any batch-affine operator, even with the true direction (\textbf{R1}; E4b fails even
oracle-directed, Result~\ref{sec:repair}) --- and this is a genuine non-invertibility, not mis-estimation, since
the estimator recovers the direction yet the operator cannot undo the per-sample noise. A \textbf{task-coupled}
nuisance subspace makes erasure destroy task signal (\textbf{R2}; the natural spatial branch,
Result~\ref{sec:rq4}, and the injected oracle-subspace erasure both fail). A shortcut that is
\textbf{under-identified from target-$X$} --- a learned reliance or a concept shift --- cannot be identified,
let alone repaired, from the marginal alone (\textbf{R3}).

\begin{table}[t]
\centering\small
\begin{tabular}{llll}
\toprule
Class & Obstruction & Invertible by target-$X$ batch-affine op? & Licensed action \\
\midrule
R0 & none (affine-in-$z$, task-sep., identified) & yes (E4; narrow, construction-matched) & \textsc{repair} \\
R1 & per-sample noise indep.\ of $z$ (2nd-moment) & no (even oracle direction) & \textsc{refuse / need per-sample info} \\
R2 & task-coupled subspace ($S\not\!\perp T$) & no (erasure destroys task) & \textsc{refuse / task-preserving inverse} \\
R3 & under-identified from target-$X$ (learned/concept) & not even identifiable & \textsc{require stronger contract} \\
\bottomrule
\multicolumn{4}{l}{\footnotesize Non-exclusive; the binding (most-restrictive) obstruction sets the action. R0 is the unique cell where none hold.}
\end{tabular}

\caption{Recoverability obstruction classes (non-exclusive; the binding obstruction sets the licensed action).
Deployable target-$X$-only batch-affine repair reaches only R0; R1/R2/R3 are detected, localized, and refused,
and require stronger information (per-sample side information, a task-preserving inverse, or labels/paired/
randomized data). ``Not repairable'' is scoped to batch-affine operators, not proven for all operators.}
\label{tab:recover}
\end{table}

\paragraph{Head-only learned-reliance bridge.}
To bridge the controlled injected shortcuts above and a full learned-reliance stress test, we ran two
CPU-only head-level probes on the frozen Phase-4B EEG latents (retraining \emph{only} a linear head;
target labels used for final scoring only). Source subject-class \emph{prevalence reweighting} on true
labels induced the subject--class correlation exactly, but did \emph{not} induce a learned subject
reliance: the task signal remained a sufficient statistic and the learnability gate failed closed. A
stronger subject-keyed, \emph{task-conflicting} relabeling (global label histogram held exact) was
\emph{learnable in-sample} --- the head fit the relabeled subset far beyond a task-only floor ($0.70$ vs
$0.20$) --- yet its held-out and target true-task harm did \emph{not} exceed a subject-scrambled control
and was \emph{not} localized to the subject subspace (it beat only a matched random-noise control, on
both datasets). Thus, in this head-only setting, neither prevalence skew nor subject-keyed task-conflict
turned naturally present subject leakage into transferable, subject-structure-specific harmful reliance;
both sit in the under-identified region (\textbf{R3}), and neither licenses the paused learned-reliance
refit. These probes do \emph{not} prove that learned subject shortcuts cannot occur; they show that under
the cheapest controlled head-level manipulations, the available natural subject signal is not readily
weaponized without stronger information or a different mechanism.

Two consequences frame the paper's contribution. First, repairability is a property of the \emph{corruption
structure $\times$ operator class $\times$ available information}, not of the shortcut's mere detectability: the
same subject signal is task-useful naturally (R2, refuse), repairable as a deterministic injected offset (R0), and
un-repairable as a stochastic injected perturbation (R1). Second, the open cases (R1/R2/R3) are exactly where FSR
\emph{refuses} or requires a stronger information contract --- the same discipline (available information $\to$
what is identifiable $\to$ what action is licensed) that recurs across leakage measurement, erasure, and
adaptation. Repair is not a switch to flip; it is an action licensed only inside a recoverability class.
